\documentclass[11pt]{article}

\usepackage[margin=1in]{geometry}
\usepackage{amsmath,amssymb,amsthm,mathtools}
\usepackage{mathrsfs}
\usepackage{booktabs}
\usepackage{array}
\usepackage{microtype}
\usepackage{hyperref}

\hypersetup{hidelinks}

\newtheorem{theorem}{Theorem}[section]
\newtheorem{proposition}[theorem]{Proposition}
\newtheorem{corollary}[theorem]{Corollary}
\newtheorem{lemma}[theorem]{Lemma}
\newtheorem{conjecture}[theorem]{Conjecture}
\theoremstyle{definition}
\newtheorem{definition}[theorem]{Definition}
\theoremstyle{remark}
\newtheorem{remark}[theorem]{Remark}

\newcommand{\cT}{\mathcal{T}}
\newcommand{\cI}{\mathcal{I}}
\newcommand{\Sym}{\operatorname{Sym}}
\newcommand{\tr}{\operatorname{tr}}
\newcommand{\Riem}{\mathcal{R}}
\newcommand{\BRiem}{\widetilde{R}}

\title{\textbf{A Curvature-Deformed Arcsine Coordinate for Cubic Geometric Refinement}\\[0.5em]
\large Exact Spherical Conjugacy, Riemannian Expansion Through Fifth Order,\\
and a High-Precision Weight-Six Reconstruction}
\author{Wayne Baker\\Independent researcher}
\date{August 8, 2026}

\begin{document}

\maketitle

\begin{abstract}
A normalized geodesic side law on a Riemannian surface defines, by inversion, a
local linearizing coordinate for geometric encode--scale--decode refinement.
In the Euclidean limit this coordinate is the ordinary arcsine function.
This paper develops the resulting \emph{curvature-deformed arcsine}
\[
\Psi_{\rho,w}(z)=\frac12 H_{\rho,w}^{-1}(2z),
\]
where \(H_{\rho,w}\) is the normalized signed side law for equal-radius geodesic
endpoints based at \(p\) and centered about a unit tangent direction \(w\).

The construction gives an exact Schr\"oder/Koenigs conjugacy
\[
\Psi_{\rho,w}\!\left(\cT_{\rho,w,q}(z)\right)
=q\,\Psi_{\rho,w}(z),
\]
and hence an exact local multiplicative semigroup.  On a sphere the deformed
arcsine is available in closed form,
\[
\Psi_\alpha(z)
=
\arcsin\!\left(
\frac{\sin(z\sin\alpha)}{\sin\alpha}
\right),
\]
which interpolates between \(\arcsin z\) in the flat limit and the identity
map at the hemispherical value \(\alpha=\pi/2\).

For a general Riemannian surface the variable-curvature expansion is derived
analytically through geometric weight five.  The coefficient functions are
finite Laurent polynomials in \(r=\sqrt{1-z^2}\), and therefore every
geometric-weight row of the Taylor coefficients is a finite combination of
generalized binomial sequences.  At weight six, where the available
geodesic-distance expansion no longer directly supplies the required
degree-eight side-law term, a high-precision intrinsic torus computation is
used to recover a rational candidate formula in eight natural curvature
invariants.  Independent geometry and angle holdouts, including a failed
restricted ansatz and its subsequent correction, provide a stringent
computational validation.  Conditional on that recovered side law, exact
symbolic inversion yields the complete weight-six Laurent kernels and
all Taylor coefficients at that weight.

The results motivate a finite Laurent-kernel conjecture for the full
curvature hierarchy and extend the cubic \(N\)-series coefficient-lifting
mechanism from the ordinary arcsine to its Riemannian deformation.
\end{abstract}

\noindent\textbf{2020 Mathematics Subject Classification.} 53B20, 53A55, 39B12, 65B05.\par
\noindent\textbf{Key words and phrases.} Riemannian normal coordinates, geodesic distance, Schr\"oder equation, Koenigs linearization, Gaussian curvature, asymptotic expansion, geometric refinement.

\section{Introduction}
\label{sec:introduction}

A recurring feature of cubic geometric refinement is that a nonlinear
geometric construction becomes linear after passing to the correct local
coordinate.  In the flat circular problem that coordinate is the arcsine:
the classical triple-angle polynomial
\[
3z-4z^3
\]
is conjugate to multiplication by \(3\) through
\[
\arcsin(3z-4z^3)=3\arcsin z
\]
on the appropriate local branch.

The same mechanism persists intrinsically on a Riemannian surface, but the
linearizing coordinate is no longer the ordinary arcsine.  The geometry
itself deforms it.

Let \(p\) be a point of an oriented Riemannian surface, \(w\in T_pM\) a unit
tangent direction, and \(\rho>0\) a sufficiently small geodesic radius.
The equal-radius signed side law
\[
c_{\rho,w}(x)
\]
records the signed intrinsic distance between radial geodesic endpoints whose
angular separation is \(x\), centered about the direction \(w\).  Normalize it
by
\[
H_{\rho,w}(x)
=
\frac{c_{\rho,w}(x)}{c'_{\rho,w}(0)}.
\]
The associated geometric encode--scale--decode map at scale \(q>0\) is
\[
\Phi_{\rho,w,q}
=
H_{\rho,w}^{-1}\circ(qH_{\rho,w}).
\]
Thus
\[
H_{\rho,w}\!\left(\Phi_{\rho,w,q}(x)\right)
=
qH_{\rho,w}(x)
\]
holds identically wherever the local inverse is defined.

This observation has two consequences.  First, the normalized side law is
already a Koenigs coordinate for the refinement family, equivalently a local
solution of Schr\"oder's equation; no separate linearization problem has to be
solved.  Second, after a half-angle
normalization the flat coordinate becomes precisely \(\arcsin z\).  This leads
to the definition
\[
\Psi_{\rho,w}(z)
=
\frac12H_{\rho,w}^{-1}(2z),
\]
which is the central object of the present paper.

The work has four layers.
\begin{enumerate}
\item The normalized side law gives an exact local semigroup and an exact
      curvature-deformed arcsine coordinate.
\item On a sphere, the coordinate and the resulting scale maps are available
      in closed form.
\item On a general Riemannian surface, the variable-curvature deformation is
      expanded analytically through geometric weight five.
\item At weight six, an independent high-precision intrinsic torus
      calculation reconstructs a rational candidate for the new side-law term at
      high precision.  Exact symbolic algebra then propagates that
      candidate to the full weight-six linearizer and to all Taylor degrees.
\end{enumerate}

The fourth item is deliberately separated in status from the first three.
The weight-six side-law coefficients have not yet been obtained from an
independent degree-eight normal-coordinate derivation.  They are therefore
stated as a high-precision computational candidate.  All algebraic consequences
drawn from that candidate are exact.

The Riemann-normal-coordinate calculations through the lower orders use the
corrected geodesic arc-length expansion developed by Brewin and its public
source repository \cite{Brewin2009,BrewinRepository}.  The abstract local
linearization mechanism belongs to the classical theory of the Schr\"oder
equation; see, for example, Iserles \cite{Iserles1993}.  Likewise, the
organization of normal-coordinate coefficients by natural curvature tensors
is classical; a recent formulation is given by Graham and Kuo
\cite{GrahamKuo2026}.  No novelty is claimed for those general frameworks.
The specific subject here is their realization by this normalized intrinsic
side law, the resulting spherical and variable-curvature linearizing
coordinate, and the finite Laurent structure of its coefficient functions.
The spherical, fifth-order curvature, and flat cubic coefficient-lifting
results build on the author's preceding notes \cite{BakerSphere,BakerCurvature,BakerNSeries}.

\section{Geometric side law and exact linearization}
\label{sec:side-law}

All statements below are local.  We take \(\rho>0\) sufficiently small that
the relevant radial endpoints and their joining short geodesics remain in a
common convex normal neighbourhood of \(p\), and all inverse maps are taken on
the local branch through the origin.

We fix the curvature convention
\[
R(A,B)C
=\nabla_A\nabla_BC-\nabla_B\nabla_AC-\nabla_{[A,B]}C,
\qquad
\Riem(A,B,C,D)=\langle R(A,B)C,D\rangle,
\]
so that on a surface
\[
\Riem(A,B,B,A)
=K\bigl(|A|^2|B|^2-\langle A,B\rangle^2\bigr).
\]
Thus, in components,
\[
\Riem_{abcd}=K(g_{ad}g_{bc}-g_{ac}g_{bd}).
\]
Brewin's displayed arc-length formulas use the opposite four-slot component
sign in the slot convention employed in the lower-order derivation; when
those formulas are invoked we therefore use
\[
\BRiem_{abcd}:=-\Riem_{abcd}
=K(g_{ac}g_{bd}-g_{ad}g_{bc}).
\]
This convention lock is the one used in the fifth-order source calculation
\cite{BakerCurvature}.

Let \((M,g)\) be a smooth oriented Riemannian surface.  Fix
\[
p\in M,\qquad |w|_g=1,
\]
and let \(J\) denote rotation by \(+\pi/2\) in the oriented tangent plane.

For sufficiently small \(x\), define
\[
u_x
=
\cos\frac{x}{2}\,w
+
\sin\frac{x}{2}\,Jw,
\qquad
v_x
=
\cos\frac{x}{2}\,w
-
\sin\frac{x}{2}\,Jw.
\]
Let
\[
P_x=\exp_p(\rho u_x),
\qquad
Q_x=\exp_p(\rho v_x).
\]
The signed side law \(c_{\rho,w}(x)\) is the signed length of the unique short
geodesic from \(P_x\) to \(Q_x\), with sign chosen so that
\[
c_{\rho,w}(-x)=-c_{\rho,w}(x).
\]
Hence \(c_{\rho,w}\) is locally odd.

Define
\[
H_{\rho,w}(x)
=
\frac{c_{\rho,w}(x)}{c'_{\rho,w}(0)}.
\]
Then
\[
H_{\rho,w}(0)=0,
\qquad
H'_{\rho,w}(0)=1,
\]
so \(H_{\rho,w}\) has a local inverse.  It is convenient to introduce
the normalized half-side coordinate
\[
\boxed{
S_{\rho,w}(u):=\frac12 H_{\rho,w}(2u).
}
\]
Then \(S_{\rho,w}(0)=0\), \(S'_{\rho,w}(0)=1\), and the
deformed arcsine introduced below is simply
\[
\Psi_{\rho,w}=S_{\rho,w}^{-1}.
\]

\begin{proposition}[Exact geometric linearization]
\label{prop:exact-linearization}
For every admissible local scale \(q>0\), define
\[
\Phi_{\rho,w,q}
=
H_{\rho,w}^{-1}\circ(qH_{\rho,w}).
\]
Then
\[
H_{\rho,w}\!\left(\Phi_{\rho,w,q}(x)\right)
=
qH_{\rho,w}(x),
\]
and the maps obey the local semigroup law
\[
\Phi_{\rho,w,q}\circ\Phi_{\rho,w,s}
=
\Phi_{\rho,w,qs}.
\]
\end{proposition}

\begin{proof}
Both statements follow directly from the definition and local invertibility
of \(H_{\rho,w}\).
\end{proof}

Writing \(q=e^t\) gives a local flow
\[
\Phi_t
=
H_{\rho,w}^{-1}(e^tH_{\rho,w}),
\]
whose infinitesimal generator is
\[
V_{\rho,w}(x)
=
\frac{H_{\rho,w}(x)}{H'_{\rho,w}(x)}.
\]
In the flat circular case,
\[
H_0(x)=2\sin\frac{x}{2},
\qquad
V_0(x)=2\tan\frac{x}{2}.
\]

\section{The curvature-deformed arcsine}
\label{sec:deformed-arcsine}

\begin{definition}
The curvature-deformed arcsine associated with \((p,w,\rho)\) is
\[
\boxed{
\Psi_{\rho,w}(z)
=
\frac12H_{\rho,w}^{-1}(2z).
}
\]
\end{definition}

Equivalently,
\[
\Psi_{\rho,w}=S_{\rho,w}^{-1},
\qquad
\Psi_{\rho,w}^{-1}(y)
=
S_{\rho,w}(y)
=
\frac12H_{\rho,w}(2y).
\]
Define
\[
\cT_{\rho,w,q}
=
\Psi_{\rho,w}^{-1}\circ(q\Psi_{\rho,w}).
\]

\begin{theorem}[Exact Schr\"oder relation]
\label{thm:schroder}
For every admissible \(q>0\),
\[
\boxed{
\Psi_{\rho,w}\!\left(\cT_{\rho,w,q}(z)\right)
=
q\,\Psi_{\rho,w}(z).
}
\]
Moreover,
\[
\cT_{\rho,w,q}\circ\cT_{\rho,w,s}
=
\cT_{\rho,w,qs}.
\]
\end{theorem}

In the flat limit,
\[
H_0(x)=2\sin\frac{x}{2},
\]
so
\[
\boxed{\Psi_0(z)=\arcsin z.}
\]
Thus the ordinary arcsine is the flat member of a geometrically defined
family.

\section{Exact spherical model}
\label{sec:sphere}

Let the surface be a sphere of radius \(R\), and write
\[
\alpha=\frac{\rho}{R}.
\]
The normalized spherical side law is an elementary consequence of spherical
trigonometry; see also the preceding spherical refinement note
\cite{BakerSphere}.  In the present setting its significance is that it gives
the linearizing coordinate in closed form.  Explicitly,
\[
\boxed{
H_\alpha(x)
=
\frac{2}{\sin\alpha}
\arcsin\!\left(
\sin\alpha\,
\sin\frac{x}{2}
\right).
}
\]

\begin{theorem}[Exact spherical deformed arcsine]
\label{thm:spherical-psi}
For sufficiently small \(z\),
\[
\boxed{
\Psi_\alpha(z)
=
\arcsin\!\left(
\frac{\sin(z\sin\alpha)}{\sin\alpha}
\right).
}
\]
\end{theorem}

The limiting cases are
\[
\Psi_0(z)=\arcsin z,
\qquad
\Psi_{\pi/2}(z)=z
\]
on the principal local branch.

If \(c=\cos^2\alpha\), then
\[
\Psi_\alpha(z)
=
z+\frac{c}{6}z^3
+\frac{c(c+8)}{120}z^5
+\frac{c(c^2+88c+136)}{5040}z^7
+O(z^9).
\]

The exact spherical tripling map is
\[
\cT_\alpha
=
\Psi_\alpha^{-1}\circ(3\Psi_\alpha),
\]
hence
\[
\Psi_\alpha(\cT_\alpha(z))=3\Psi_\alpha(z).
\]
Its Taylor expansion begins
\[
\cT_\alpha(z)
=
3z-4cz^3
+16c(c-1)z^5
-\frac4{15}c(c-1)(311c-221)z^7
+O(z^9).
\]
At \(c=1\) this reduces to \(3z-4z^3\), while at \(c=0\) it reduces to
\(3z\).

\section{General Riemannian expansion through weight five}
\label{sec:weight-five}

Write
\[
K=K(p),\qquad
K_w=(\nabla_wK)(p),\qquad
K_{ww}=\nabla^2K(w,w)(p),
\]
\[
D=(\Delta K)(p),\qquad
D_w=(\nabla_w\Delta K)(p),\qquad
T_{www}=\nabla^3K(w,w,w)(p).
\]
Let \(r=\sqrt{1-z^2}\).  Then
\[
\Psi_{\rho,w}(z)
=
\arcsin z
+\rho^2\Psi_2(z)
+\rho^3\Psi_3(z)
+\rho^4\Psi_4(z)
+\rho^5\Psi_5(z)
+O(\rho^6).
\]

\subsection{Analytic side-law input and inversion}

For completeness, we record the exact lower-order side-law input needed for
the inversion.  Put
\[
s=\sin u,\qquad c=\cos u.
\]
The signed side law through fifth radial order, derived from the corrected
normal-coordinate distance expansion in \cite{Brewin2009} with the convention
lock above, is \cite{BakerCurvature}
\[
c_{\rho,w}(2u)
=2\rho s\,F_{\rho,w}(u),
\]
where
\[
F_{\rho,w}(u)
=1+f_2\rho^2+f_3\rho^3+B_4\rho^4+B_5\rho^5+O(\rho^6),
\]
with
\[
f_2=-\frac16Kc^2,
\qquad
f_3=-\frac1{12}K_wc^3,
\]
\[
B_4
=\frac{c^2}{120}
\left[(1-9s^2)K^2-3c^2K_{ww}-s^2(D-K_{ww})\right],
\]
and, writing
\[
T_\Sigma:=3D_w-3T_{www}+2KK_w,
\]
\[
B_5
=\frac{c^3}{1080}
\left[(9c^2-92s^2)KK_w-6c^2T_{www}-2s^2T_\Sigma\right].
\]
Since \(c'_{\rho,w}(0)=\rho F_{\rho,w}(0)\), the normalized
half-side coordinate is
\[
\boxed{
S_{\rho,w}(u)
=s\,\frac{F_{\rho,w}(u)}{F_{\rho,w}(0)}.
}
\]
Thus the calculation of \(\Psi_{\rho,w}\) through weight five is a
pure perturbative inversion of this displayed scalar series.

\begin{theorem}[Riemannian deformed arcsine through weight five]
\label{thm:psi-five}
The coefficient functions are
\[
\boxed{
\Psi_2(z)
=-\frac{Kz^3}{6r}
=-\frac{Kz}{6}(r^{-1}-r),
}
\]
\[
\boxed{
\Psi_3(z)
=-\frac{K_wz(1-r^3)}{12r}
=-\frac{K_wz}{12}(r^{-1}-r^2),
}
\]
\[
\boxed{
\Psi_4(z)
=
\frac{z^3}{360r^3}
\left[
K^2(2r^4+13r^2+5)
+3Dr^4
-3K_{ww}r^2(4r^2+3)
\right],
}
\]
and
\[
\boxed{
\begin{aligned}
\Psi_5(z)
=
\frac{z}{360r^3}\Big[
&2D_wr^5z^2
\\
&+KK_w(1-r)
(5r^6+5r^5+8r^4+18r^3+8r^2+5r+5)
\\
&-2T_{www}r^2(1-r)
(2r^4+2r^3+r^2+r+1)
\Big].
\end{aligned}}
\]
\end{theorem}

\begin{proof}
Write
\[
u=\arcsin z+\rho^2U_2+\rho^3U_3+\rho^4U_4+\rho^5U_5+O(\rho^6)
\]
and substitute into \(S_{\rho,w}(u)=z\).  At each radial order the
unknown \(U_m\) occurs linearly with coefficient
\(\cos(\arcsin z)=r\neq0\) on the local branch.  Solving successively
for \(U_2,U_3,U_4,U_5\) and simplifying with \(z^2=1-r^2\) gives the
four displayed coefficient functions.  The calculation uses exact rational
algebra; the accompanying symbolic audit reconstructs the normalized side law
from the formulas above and verifies all four identities identically.
\end{proof}

Equivalently,
\[
\begin{aligned}
\Psi_4(z)
={}&
\frac{K^2z}{360}(-2r^3-11r+8r^{-1}+5r^{-3})
\\
&+\frac{Dz}{120}(r-r^3)
+\frac{K_{ww}z}{120}(4r^3-r-3r^{-1}),
\end{aligned}
\]
and
\[
\begin{aligned}
\Psi_5(z)
={}&
\frac{KK_wz}{360}
(-5r^4-3r^2-10r+10+3r^{-1}+5r^{-3})
\\
&+\frac{D_wz}{180}(r^2-r^4)
+\frac{T_{www}z}{360}(4r^4-2r^2-2r^{-1}).
\end{aligned}
\]

The \(D_w\) row terminates:
\[
\frac{D_wz}{180}(r^2-r^4)
=
\frac{D_w}{180}(z^3-z^5),
\]
so \(D_w\) contributes only to \(\gamma_3\) and \(\gamma_5\).

\subsection{First Taylor coefficients}

Write
\[
\Psi_{\rho,w}(z)
=
z+\gamma_3z^3+\gamma_5z^5+\gamma_7z^7+\cdots.
\]
Through weight five,
\[
\begin{aligned}
\gamma_3={}&
\frac16-\frac16K\rho^2-\frac18K_w\rho^3
\\
&+\left(\frac1{18}K^2+\frac1{120}D-\frac7{120}K_{ww}\right)\rho^4
\\
&+\left(\frac3{40}KK_w+\frac1{180}D_w-\frac7{360}T_{www}\right)\rho^5
+O(\rho^6),
\end{aligned}
\]
\[
\begin{aligned}
\gamma_5={}&
\frac3{40}-\frac1{12}K\rho^2-\frac1{32}K_w\rho^3
\\
&+\left(\frac{13}{360}K^2-\frac1{240}D+\frac1{240}K_{ww}\right)\rho^4
\\
&+\left(\frac3{160}KK_w-\frac1{180}D_w+\frac{13}{1440}T_{www}\right)\rho^5
+O(\rho^6),
\end{aligned}
\]
and
\[
\begin{aligned}
\gamma_7={}&
\frac5{112}-\frac1{16}K\rho^2-\frac5{192}K_w\rho^3
\\
&+\left(\frac7{180}K^2-\frac1{960}D-\frac1{192}K_{ww}\right)\rho^4
\\
&+\left(\frac5{144}KK_w-\frac1{576}T_{www}\right)\rho^5
+O(\rho^6).
\end{aligned}
\]

\section{Curved tripling and cubic coefficient lifting}
\label{sec:lifting}

For the cubic scale set
\[
\cT_{\rho,w}
=
\cT_{\rho,w,3}
=
\Psi_{\rho,w}^{-1}\circ(3\Psi_{\rho,w}).
\]
Writing
\[
\cT_{\rho,w}(z)
=
3z+\tau_3z^3+\tau_5z^5+\tau_7z^7+\cdots,
\]
formal conjugacy gives
\[
\boxed{\tau_3=-24\gamma_3,}
\]
\[
\boxed{\tau_5=24(27\gamma_3^2-10\gamma_5),}
\]
and
\[
\boxed{
\tau_7
=-24(945\gamma_3^3-675\gamma_3\gamma_5+91\gamma_7).
}
\]
The flat specialization gives \(\cT_0(z)=3z-4z^3\) exactly.

For \(j\ge1\), define
\[
\mathscr L_j[P](z)
=
\frac{3^{2j+1}P(z)-P(\cT_{\rho,w}(z))}
{3(3^{2j}-1)}.
\]

\begin{theorem}[Curved cubic coefficient lifting]
\label{thm:curved-lifting}
Let \(P_j\) be odd and suppose
\[
\Psi_{\rho,w}(z)-P_j(z)=O(z^{2j+1}).
\]
Then
\[
P_{j+1}=\mathscr L_j[P_j]
\]
satisfies
\[
\Psi_{\rho,w}(z)-P_{j+1}(z)=O(z^{2j+3}).
\]
\end{theorem}

\begin{proof}
Set \(E_j=\Psi_{\rho,w}-P_j\).  By the exact Schr\"oder relation,
\[
\Psi_{\rho,w}(\cT_{\rho,w}(z))=3\Psi_{\rho,w}(z).
\]
If
\[
E_j(z)=az^{2j+1}+O(z^{2j+3}),
\]
then, because \(\cT_{\rho,w}(z)=3z+O(z^3)\),
\[
E_j(\cT_{\rho,w}(z))
=a(3z)^{2j+1}+O(z^{2j+3}).
\]
The defining linear combination cancels the leading error, and oddness skips
the intervening even degree.
\end{proof}

Thus the flat cubic \(N\)-series reconstructs the ordinary arcsine coefficient
sequence, while the Riemannian cubic \(N\)-series reconstructs the
curvature-deformed sequence.

\section{Weight-six extension}
\label{sec:weight-six}

Define
\[
S_4:=\Sym\nabla^4K,
\]
with contractions
\[
S_{wwww}:=S_4(w,w,w,w),
\qquad
(\tr S)_{ww}:=(\tr S_4)(w,w),
\qquad
\tr^2S:=\tr(\tr S_4).
\]
The high-precision recovery is organized in the following natural
weight-six family:
\[
K^3,\quad K\Delta K,\quad KK_{ww},\quad |\nabla K|^2,\quad K_w^2,
\]
\[
S_{wwww},\quad (\tr S)_{ww},\quad \tr^2S.
\]
This family is motivated by geometric weight, direction/orientation symmetry,
and the classical natural-tensor framework \cite{GrahamKuo2026}.  At this
stage it is a candidate basis for this side law rather than an abstract
completeness theorem for all possible weight-six natural scalars.  The
full-rank recovery design and unused holdouts test the basis empirically.

\subsection{High-precision side-law reconstruction}

Let \(z=\sin u\) and \(r=\cos u\).  Write
\[
S_{\rho,w}(u)
=
\sin u+\rho^2s_2(u)+\cdots+\rho^6s_6(u)+O(\rho^7).
\]

\begin{proposition}[High-precision weight-six reconstruction]
\label{prop:s6}
High-precision intrinsic torus computations, rational reconstruction, and
independent geometry and angle holdout tests identify
\[
\begin{aligned}
s_6
=
z\Bigg\{&
K^3(1-r^2)
\left(\frac{31}{15120}-\frac{11}{280}r^2+\frac5{112}r^4\right)
\\
&+K\Delta K(1-r^2)
\left(-\frac1{112}r^2+\frac{127}{15120}r^4\right)
\\
&+KK_{ww}(1-r^2)
\left(\frac{29}{5040}+\frac{53}{5040}r^2-\frac{53}{1260}r^4\right)
\\
&+|\nabla K|^2(1-r^2)
\left(-\frac{17}{10080}r^2+\frac{71}{30240}r^4\right)
\\
&+K_w^2
\left(
\frac5{1008}+\frac{17}{10080}r^2-\frac1{144}r^3
-\frac{157}{5040}r^4+\frac{317}{10080}r^6
\right)
\\
&+S_{wwww}(1-r^2)
\left(\frac1{1008}+\frac1{1260}r^2+\frac1{315}r^4\right)
\\
&+(\tr S)_{ww}(1-r^2)
\left(\frac1{2520}r^2-\frac1{420}r^4\right)
\\
&+\tr^2S(1-r^2)
\left(-\frac1{5040}r^2+\frac1{5040}r^4\right)
\Bigg\}.
\end{aligned}
\]
\end{proposition}

\begin{remark}[Status]
Proposition~\ref{prop:s6} is a high-precision rational reconstruction
supported by independent holdout computations; it is not yet an analytic
degree-eight normal-coordinate theorem or an interval-arithmetic proof.  The
distinction is maintained throughout this paper.
\end{remark}

The \(K^3\) row is independently fixed by the exact spherical model:
\[
\boxed{
[s_6]_{K^3}
=
\frac{K^3z^3}{15120}(675r^4-594r^2+31).
}
\]

The \(K_w^2\) row deserves separate attention.  A first restricted ansatz in
which every row factored as
\[
z(1-r^2)(A+Br^2+Cr^4)
\]
failed at new side angles.  The discrepancy was traced to normalization of
the already-known weight-three side-law term
\[
f_3(u)=-\frac{K_w}{12}r^3.
\]
Indeed,
\[
f_3(0)^2-f_3(u)f_3(0)
=
\frac{K_w^2}{144}(1-r^3),
\]
which is not divisible by \(1-r^2\) with an even-polynomial quotient.  Once
this forced nonlinear contribution is included, the independent angle
holdouts return to the numerical extraction scale.

\subsection{Universal inverse formula at weight six}

Let
\[
S_{\rho,w}(u)
=
s_0(u)+\rho^2s_2(u)+\cdots+\rho^6s_6(u)+O(\rho^7),
\qquad s_0(u)=\sin u,
\]
and
\[
\Psi_{\rho,w}(z)
=
u_0+\rho^2U_2+\cdots+\rho^6U_6+O(\rho^7),
\qquad u_0=\arcsin z.
\]
Exact formal inversion gives
\[
\begin{aligned}
U_6=-\frac1r\Bigg[
&s_6+s_4'U_2+s_3'U_3+s_2'U_4+\frac12s_2''U_2^2
\\
&-z\left(U_2U_4+\frac12U_3^2\right)-\frac r6U_2^3
\Bigg],
\end{aligned}
\]
where side-law derivatives are evaluated at \(u=u_0\).  In particular,
\[
\boxed{s_5\ \text{and}\ U_5\ \text{do not enter}\ U_6.}
\]
Thus
\[
\Psi_6=\Psi_6^{\mathrm{ind}}-\frac{s_6}{r},
\]
where the induced term is determined entirely by weights \(2,3,4\).

\subsection{Complete finite Laurent layer}

\begin{theorem}[Conditional weight-six Laurent expansion]
\label{thm:weight-six-laurent}
Assume Proposition~\ref{prop:s6}.  Then
\[
\boxed{
\Psi_6(z;p,w)
=
z\sum_I\cI_I(p,w)L_I(r),
\qquad r=\sqrt{1-z^2},
}
\]
with
\[
L_{K^3}
=
\frac{r^5}{2835}+\frac{r^3}{5670}+\frac{23r}{15120}
+\frac1{9072r}+\frac1{6480r^3}-\frac1{432r^5},
\]
\[
L_{K\Delta K}
=
-\frac{r^5}{756}+\frac{53r^3}{15120}-\frac r{280}+\frac1{720r},
\]
\[
L_{KK_{ww}}
=
-\frac{r^5}{315}-\frac{5r^3}{504}+\frac{13r}{1680}
+\frac1{840r}+\frac1{240r^3},
\]
\[
L_{|\nabla K|^2}
=
\frac{71r^5}{30240}-\frac{61r^3}{15120}+\frac{17r}{10080},
\]
\[
\begin{aligned}
L_{K_w^2}
={}&
-\frac{r^5}{140}+\frac{139r^3}{10080}-\frac{r^2}{48}
-\frac{17r}{10080}
\\
&+\frac1{72}-\frac1{672r}+\frac1{288r^3},
\end{aligned}
\]
\[
L_{S_{wwww}}
=
\frac{r^5}{315}-\frac{r^3}{420}+\frac r{5040}-\frac1{1008r},
\]
\[
L_{(\tr S)_{ww}}
=
-\frac{r^5}{420}+\frac{r^3}{360}-\frac r{2520},
\]
and
\[
L_{\tr^2S}
=
\frac{r^5}{5040}-\frac{r^3}{2520}+\frac r{5040}.
\]
Each kernel obeys \(L_I(1)=0\).
\end{theorem}

\begin{proof}
Using Proposition~\ref{prop:s6}, substitute the previously established
\(U_2,U_3,U_4\) into the universal weight-six inverse formula.  Exact symbolic
simplification gives the displayed kernels.  This inversion step uses no
numerical approximation.
\end{proof}

\subsection{Exact spherical checksum at weight six}

The constant-curvature sector of Theorem~\ref{thm:weight-six-laurent} can be
checked independently from the exact spherical formula, without using the
torus recovery.  Expanding
\[
\Psi_\alpha(z)
=
\arcsin\!\left(
\frac{\sin(z\sin\alpha)}{\sin\alpha}
\right)
\]
through order \(\alpha^6\) gives
\[
\boxed{
[\alpha^6]\Psi_\alpha(z)
=
z\left(
\frac{r^5}{2835}
+\frac{r^3}{5670}
+\frac{23r}{15120}
+\frac{1}{9072r}
+\frac{1}{6480r^3}
-\frac{1}{432r^5}
\right),
}
\qquad r=\sqrt{1-z^2}.
\]
Thus
\[
[\alpha^6]\Psi_\alpha(z)=zL_{K^3}(r)
\]
identically as a function of \(z\).  Since on a constant-curvature surface
\[
\alpha=\rho\sqrt K,
\qquad
\alpha^6=K^3\rho^6,
\]
the entire \(K^3\rho^6\) row in the general weight-six Laurent expansion is
fixed independently by the exact spherical model.  Exact symbolic
simplification of the difference between the two expressions gives zero.

\section{All Taylor degrees at weight six}
\label{sec:all-n}

Define
\[
\boxed{
\beta_n^{(\nu)}=(-1)^n\binom{\nu/2}{n},
}
\]
so that
\[
(1-z^2)^{\nu/2}
=
\sum_{n\ge0}\beta_n^{(\nu)}z^{2n}.
\]

\begin{corollary}[All-\(n\) weight-six coefficient rows]
\label{cor:all-n}
Assume Proposition~\ref{prop:s6}.  For every \(n\ge1\),
\[
\begin{aligned}
[\rho^6]\gamma_{2n+1}
={}&
A_nK^3+B_nK\Delta K+C_nKK_{ww}+D_n|\nabla K|^2
\\
&+E_nK_w^2+F_nS_{wwww}+G_n(\tr S)_{ww}+H_n\tr^2S,
\end{aligned}
\]
where
\[
\begin{aligned}
A_n={}&
\frac1{2835}\beta_n^{(5)}
+\frac1{5670}\beta_n^{(3)}
+\frac{23}{15120}\beta_n^{(1)}
\\
&+\frac1{9072}\beta_n^{(-1)}
+\frac1{6480}\beta_n^{(-3)}
-\frac1{432}\beta_n^{(-5)},
\end{aligned}
\]
\[
B_n
=
-\frac1{756}\beta_n^{(5)}
+\frac{53}{15120}\beta_n^{(3)}
-\frac1{280}\beta_n^{(1)}
+\frac1{720}\beta_n^{(-1)},
\]
\[
\begin{aligned}
C_n={}&
-\frac1{315}\beta_n^{(5)}
-\frac5{504}\beta_n^{(3)}
+\frac{13}{1680}\beta_n^{(1)}
\\
&+\frac1{840}\beta_n^{(-1)}
+\frac1{240}\beta_n^{(-3)},
\end{aligned}
\]
\[
D_n
=
\frac{71}{30240}\beta_n^{(5)}
-\frac{61}{15120}\beta_n^{(3)}
+\frac{17}{10080}\beta_n^{(1)},
\]
\[
\begin{aligned}
E_n={}&
-\frac1{140}\beta_n^{(5)}
+\frac{139}{10080}\beta_n^{(3)}
-\frac1{48}\beta_n^{(2)}
-\frac{17}{10080}\beta_n^{(1)}
\\
&+\frac1{72}\beta_n^{(0)}
-\frac1{672}\beta_n^{(-1)}
+\frac1{288}\beta_n^{(-3)},
\end{aligned}
\]
\[
F_n
=
\frac1{315}\beta_n^{(5)}
-\frac1{420}\beta_n^{(3)}
+\frac1{5040}\beta_n^{(1)}
-\frac1{1008}\beta_n^{(-1)},
\]
\[
G_n
=
-\frac1{420}\beta_n^{(5)}
+\frac1{360}\beta_n^{(3)}
-\frac1{2520}\beta_n^{(1)},
\]
and
\[
H_n
=
\frac1{5040}\beta_n^{(5)}
-\frac1{2520}\beta_n^{(3)}
+\frac1{5040}\beta_n^{(1)}.
\]
\end{corollary}

\begin{proof}
Every kernel in Theorem~\ref{thm:weight-six-laurent} is a finite linear
combination of powers \(r^\nu=(1-z^2)^{\nu/2}\).  Expanding each power by the
generalized binomial theorem and collecting \(z^{2n+1}\) gives the result.
\end{proof}

\subsection{Low Taylor degrees through weight six}

Combining the analytic lower-order results with the weight-six extension gives
\[
\begin{aligned}
\gamma_3={}&
\frac16-\frac16K\rho^2-\frac18K_w\rho^3
\\
&+\left(\frac1{18}K^2+\frac1{120}D-\frac7{120}K_{ww}\right)\rho^4
\\
&+\left(\frac3{40}KK_w+\frac1{180}D_w-\frac7{360}T_{www}\right)\rho^5
\\
&+\Bigg(
-\frac1{135}K^3+\frac1{1890}K\Delta K+\frac{13}{504}KK_{ww}
-\frac1{1512}|\nabla K|^2
\\
&\hspace{1.4cm}
+\frac{47}{2016}K_w^2-\frac5{1008}S_{wwww}
+\frac1{504}(\tr S)_{ww}
\Bigg)\rho^6+O(\rho^7),
\end{aligned}
\]
\[
\begin{aligned}
\gamma_5={}&
\frac3{40}-\frac1{12}K\rho^2-\frac1{32}K_w\rho^3
\\
&+\left(\frac{13}{360}K^2-\frac1{240}D+\frac1{240}K_{ww}\right)\rho^4
\\
&+\left(\frac3{160}KK_w-\frac1{180}D_w+\frac{13}{1440}T_{www}\right)\rho^5
\\
&+\Bigg(
-\frac1{108}K^3-\frac1{5040}K\Delta K-\frac1{420}KK_{ww}
+\frac3{1120}|\nabla K|^2
\\
&\hspace{1.4cm}
-\frac{83}{40320}K_w^2+\frac{47}{10080}S_{wwww}
-\frac{17}{5040}(\tr S)_{ww}+\frac1{5040}\tr^2S
\Bigg)\rho^6+O(\rho^7),
\end{aligned}
\]
and
\[
\begin{aligned}
\gamma_7={}&
\frac5{112}-\frac1{16}K\rho^2-\frac5{192}K_w\rho^3
\\
&+\left(\frac7{180}K^2-\frac1{960}D-\frac1{192}K_{ww}\right)\rho^4
\\
&+\left(\frac5{144}KK_w-\frac1{576}T_{www}\right)\rho^5
\\
&+\Bigg(
-\frac{227}{15120}K^3+\frac{13}{10080}K\Delta K+\frac3{320}KK_{ww}
-\frac{11}{10080}|\nabla K|^2
\\
&\hspace{1.4cm}
+\frac{119}{11520}K_w^2-\frac{59}{40320}S_{wwww}
+\frac{19}{20160}(\tr S)_{ww}-\frac1{10080}\tr^2S
\Bigg)\rho^6+O(\rho^7).
\end{aligned}
\]

At weight five,
\[
\frac{D_wz}{180}(r^2-r^4)=\frac{D_w}{180}(z^3-z^5),
\]
so \([D_w]\gamma_{2n+1}=0\) for \(n\ge3\).

At weight six, \(\beta_n^{(0)}=0\) for \(n\ge1\), so the constant term in
\(L_{K_w^2}\) affects only the linear normalization.  Also
\(\beta_n^{(2)}=0\) for \(n\ge2\), so the term \(-r^2/48\) contributes only
to \(\gamma_3\).  A further cancellation is
\[
[\tr^2S\,\rho^6]\gamma_3=0.
\]

\section{Computational validation of the weight-six layer}
\label{sec:verification}

The weight-six side-law candidate was recovered independently of the formal
normal-coordinate inversion formulas using the ring-torus metric
\[
ds^2=a^2\,d\theta^2+(R+a\cos\theta)^2\,d\phi^2.
\]
Intrinsic radial endpoints and joining geodesics were computed directly at
arbitrary precision.

The normalization derivative was obtained from the scalar Jacobi equation
along the central radial geodesic,
\[
j''+Kj=0,
\qquad j(0)=0,
\qquad j'(0)=1,
\]
so that
\[
c'_{\rho,w}(0)=j(\rho).
\]
Opposite directions \(w\) and \(-w\) were averaged, cancelling odd geometric
weights.  After subtraction of the known weight-two and weight-four terms,
\[
E_6(\rho)
=
\frac{S_{\mathrm{even}}-\sin u-\rho^2s_2-\rho^4s_4}{\rho^6}
=
s_6+O(\rho^2).
\]

A higher even extrapolation basis
\[
1,\rho^2,\rho^4,\rho^6,\rho^8
\]
was used to recover \(s_6\).  Ten torus geometries at three initial side
angles produced thirty measurements.  Twenty-one equations from the first
seven geometries determined the unknown coefficients, while nine equations
from three unused geometries served as holdouts.  After higher-order
reprocessing, the holdout root-mean-square residual was approximately
\[
9.6\times10^{-23}.
\]

A second validation used side angles not present in the recovery.  The
initial restricted angular ansatz failed by many orders of magnitude beyond
the extraction error.  The discrepancy was proportional to \(K_w^2\) across
all tested geometries and was traced to the normalization term
\[
\frac{K_w^2}{144}(1-r^3).
\]
After that mechanism was incorporated, twelve independent angle holdouts
collapsed to a root-mean-square residual of approximately
\[
3.5\times10^{-24}.
\]

Finally, exact symbolic arithmetic verified the normalization
\[
L_I(1)=0,
\]
the exact spherical \(K^3\) kernel, and the all-\(n\) formulas against direct
Taylor expansion through \(n=12\).

The accompanying reproducibility archive contains separate scripts for the
weight-six inverse scaffold, angular-target and rank audits, high-precision
torus pilot, full recovery, higher-order reprocessing, out-of-sample angular
validation, direct \(K_w^2\) kernel isolation, corrected-kernel check, and
exact Laurent/all-\(n\) audit, together with the raw and processed data used
for the reported holdouts.

\section{Finite Laurent-kernel structure}
\label{sec:finite-laurent}

Through weights two to six the deformed arcsine has the observed structure
\[
\boxed{
\Psi_m(z;p,w)
=
z\sum_I\cI_{m,I}(p,w)
L_{m,I}\!\left(\sqrt{1-z^2}\right),
}
\]
where \(\cI_{m,I}\) is a natural curvature invariant of geometric weight
\(m\), and every coefficient function encountered so far is a finite Laurent
polynomial.

The weight-six calculation is a stronger test than the lower orders because
a more restrictive angular hypothesis was actually falsified.  The failure
was localized to a single invariant sector, explained by a nonlinear
normalization effect already implied by the lower-order geometry, and
corrected without changing the invariant family.  Independent holdouts then
returned to the numerical extraction scale.

\begin{conjecture}[Finite Laurent-kernel conjecture]
\label{conj:finite-laurent}
At every finite geometric weight \(m\),
\[
\boxed{
\Psi_m(z;p,w)
=
z\sum_I\cI_{m,I}(p,w)L_{m,I}(r),
\qquad r=\sqrt{1-z^2},
}
\]
where each \(\cI_{m,I}\) is an admissible natural curvature invariant of
weight \(m\), and each \(L_{m,I}\) is a finite Laurent polynomial in \(r\).
\end{conjecture}

If the conjecture holds, then every geometric-weight row of
\[
\gamma_3,\gamma_5,\gamma_7,\ldots
\]
is a finite linear combination of generalized binomial sequences
\[
\beta_n^{(\nu)}=(-1)^n\binom{\nu/2}{n}.
\]
Finite-support selection rules arise whenever a nonnegative integral exponent
\(\nu\) causes the associated generalized binomial sequence to terminate.

\section{Discussion}
\label{sec:discussion}

The central point of this paper is not that local analytic linearization is
new.  Schr\"oder and Koenigs linearization belong to classical local dynamics.
What is geometrically significant here is that the normalized intrinsic side
law supplies the linearizing coordinate directly:
\[
H_{\rho,w}(\Phi_{\rho,w,q}(x))=qH_{\rho,w}(x).
\]
The corresponding half-angle inverse is a curvature-deformed arcsine.

The exact sphere provides a closed model in which curvature continuously
deforms the flat arcsine toward the identity map.  The general Riemannian
expansion shows how the local curvature jet enters this coordinate.  Through
weight five the result is analytic.  At weight six the intrinsic
high-precision calculation reconstructs a candidate invariant layer whose
subsequent Laurent and all-\(n\) consequences are exact conditional on that
input.

The strongest structural observation is the persistence of finite Laurent
kernels.  This turns a local geometric expansion into a statement about an
entire infinite family of Taylor coefficients at once.  It also explains why
selection rules can occur: finite nonnegative powers of \(r\) correspond to
terminating generalized-binomial sequences.

The remaining major proof obligation is the independent derivation of the
degree-eight geodesic-distance term needed to establish
Proposition~\ref{prop:s6} analytically.  Such a derivation would convert the
present high-precision weight-six reconstruction into a theorem and would
provide the first direct analytic test of
Conjecture~\ref{conj:finite-laurent} beyond the orders from which the pattern
was discovered.

\section{Conclusion}
\label{sec:conclusion}

The normalized geodesic side law defines an exact local linearizing coordinate
for geometric refinement on a Riemannian surface.  After half-angle
normalization this coordinate becomes
\[
\Psi_{\rho,w}(z)=\frac12H_{\rho,w}^{-1}(2z),
\]
which reduces to the ordinary arcsine in the flat limit.

The resulting framework gives:
\begin{itemize}
\item an exact local multiplicative semigroup;
\item an exact spherical curvature-deformed arcsine;
\item an analytic variable-curvature expansion through weight five;
\item a curved cubic coefficient-lifting theorem;
\item a high-precision reconstruction of the weight-six side-law candidate;
\item exact weight-six finite Laurent kernels conditional on that candidate;
\item all Taylor coefficients at weight six in generalized-binomial form; and
\item a sharpened finite Laurent-kernel conjecture for the full curvature
      hierarchy.
\end{itemize}

The exact spherical solution, variable-curvature expansion,
coefficient-lifting mechanism, and weight-six computational validation all
point to a common interpretation: cubic geometric refinement does not lose
the flat arcsine linearizer under curvature.  Rather, curvature deforms that
linearizer in a structured invariant hierarchy.

\begin{thebibliography}{99}

\bibitem{Brewin2009}
L.~C. Brewin,
\newblock ``Riemann normal coordinate expansions using Cadabra,''
\newblock \emph{Classical and Quantum Gravity} \textbf{26} (2009), 175017;
corrected arXiv version arXiv:0903.2087v4 (2022).
\newblock doi:10.1088/0264-9381/26/17/175017.

\bibitem{BrewinRepository}
L.~C. Brewin,
\newblock \emph{riemann-normal-coords: LaTeX and Cadabra sources for Riemann normal coordinate expansions},
\newblock public GitHub repository, accessed August 8, 2026.
\newblock \url{https://github.com/leo-brewin/riemann-normal-coords}.

\bibitem{Iserles1993}
A.~Iserles,
\newblock ``A constructive approach to the Schr\"oder equation,''
\newblock \emph{Journal of Computational and Applied Mathematics}
\textbf{46} (1993), 301--314.
\newblock doi:10.1016/0377-0427(93)90304-T.

\bibitem{GrahamKuo2026}
C.~R. Graham and T.-M. Kuo,
\newblock ``Geodesic normal coordinates and natural tensors for pseudo-Riemannian submanifolds,''
\newblock \emph{Proceedings of the American Mathematical Society}
\textbf{154} (2026), 339--351.
\newblock doi:10.1090/proc/17432.

\bibitem{BakerSphere}
W.~Baker,
\newblock ``A Spherical Extension of the Baker Cubic Trisection Refinement Map,''
\newblock unpublished manuscript, August 2026.

\bibitem{BakerCurvature}
W.~Baker,
\newblock ``Directional Curvature Expansion of the Baker Cubic Coefficient on Riemannian Surfaces Through Fifth Order,''
\newblock unpublished manuscript, August 2026.

\bibitem{BakerNSeries}
W.~Baker,
\newblock ``Cubic \(N\)-Series Acceleration and Newton-Coefficient Lifting for the Arcsine,''
\newblock unpublished manuscript, 2026.

\end{thebibliography}

\end{document}
